\section{Introduction and problem setting}\label{sec:foundations}

Many leader--follower models contain a large number of local follower
choices but only a few shared resources, aggregate performance measures,
and leader decisions. A tariff may influence the operation of many units,
while their nonlinear interaction depends only on total load or emissions.
The relevant question is whether the number of shared quantities can
control the complexity of global bilevel optimization even as the number
of follower variables grows. The distinction between the dimension of the
follower and the dimension needed to describe its responses is central.

A tractable follower problem alone does not answer this question. The
leader optimizes over the responses induced by all its possible decisions;
this optimization can be nonconvex even if every follower has a unique
convex minimizer. Conversely, nonconvexity confined to a fixed number of
aggregate coordinates need not prevent an exact description of all global
responses. The development below separates three tasks: describing global
follower responses, optimizing over those responses under a specified
selection convention, and representing the requested output at a stated
accuracy.

A local variable can represent a unit's load, production or resource use;
a small block can represent several jointly constrained choices of one unit.
The shared quantities can include total load and a few resource balances.
We model one follower optimizing the joint allocation. Separate local agents
are also represented when their decisions decouple, but a nonconvex aggregate
term does not in general describe a Nash equilibrium among independent agents.
The mathematical results concern the stated continuous models; the examples
motivate their structure without asserting an application-specific calibration.

\subsection{Contributions and scope}\label{subsec:contributions}

The main result is that a large follower dimension can be replaced by a fixed
\emph{response description dimension} under explicit local structure.
Our original contributions are the combined complexity and recovery guarantees
below, together with the restriction-preserving boundary constructions.
Their relationship to established ingredients is detailed in
\cref{subsec:related} and Table~\ref{tab:prior-comparison}.

\paragraph{Exact global responses.}
For positive-definite quadratic local blocks of fixed dimension, few shared
linear resource rows, and a polynomial cost of few aggregates, we eliminate
all local variables into rational formulas. The number of blocks and local
rows can grow. The aggregate cost can be nonconvex. The decisive bound is a
polynomial number of realizable joint regimes in fixed leader, aggregate and
multiplier dimension; independent enumeration of all local active sets would
not give this bound. Global value comparisons within that compressed space
retain all follower optima, not just stationary points.
\Cref{thm:exact-block,thm:exact-semantics} give exact upper optimization,
selection and attainment decisions, and polynomial-size common-field output.
Moving normals preserve compression, while fixed normals give optimistic
attainment. With fixed constraint normals, constant local Hessians and fixed input degree,
the degree of a selected output field is independent of the number of follower variables
(\cref{cor:constant-hessian-degree}); \cref{cor:robust-degree} extends this
arithmetic refinement to the stated moving shared and measurement normals.
\Cref{app:fixed-core} gives a related polyhedral elimination theorem and a
constructive recovery by a common convex combination of a few support tuples.

\paragraph{Robustness and certified use of approximate structure.}
Cost-near-optimal followers need not be stationary. We instead minimize on
fibers of the measurements used by each upper criterion. This gives exact
robust feasibility and worst-objective optimization with a fixed measurement
count \emph{per criterion}, even if their combined rank grows
(\cref{thm:nearoptimal-robust}). A separate perturbation argument certifies
coordinate statuses of a dense convex quadratic response over whole leader
cells. If a cover has $M$ cells and at most $t$ uncertified coordinates per
cell, exact recovery needs at most $M3^t$ LPs after screening
(\cref{thm:screening-recovery}). A computable neighborhood bounds $t$ under a
transition-multiplicity condition; small matrix error alone is insufficient.

\paragraph{Global approximation in accuracy bits.}
For strictly convex polynomial local costs, affine resource data, and a convex
polynomial aggregate, \cref{thm:accuracy-global} returns a rational leader
with a global additive guarantee for an arbitrary explicitly encoded polynomial
upper objective. The running time is polynomial in input length, numerical
degree and $B$ for error $2^{-B}$. No uniform positive curvature or Slater
condition is assumed. The proof supplies uniform rational inverse
approximations, residual-to-response bounds, and rational recovery that
preserves the base feasible leader polytope. An effective response modulus has
the sharp exponent $1/P$, where $P$ is the largest local marginal degree
(\cref{prop:accuracy-response-modulus}). Response-dependent upper inequalities
admit outer and inner guarantees; exact feasible near-optimal rational output
requires the stated tightening, convex-anchor or reserve conditions.

\paragraph{Structural and arithmetic boundaries.}
Exact optimization is NP-hard already with one leader, a unit-box follower,
an affine upper objective and a dense positive-definite Hessian arbitrarily
close to the identity (\cref{thm:conditioned-hardness}). A complementary
condition-dependent scheme has polynomial dependence on inverse normalized
additive error, which is consistent with that exact hardness
(\cref{thm:conditioned-grid}). We also give bounded-core positive results,
distinguish paths in leader interactions from paths in follower constraints,
and prove a padded identity-Hessian path construction with a fixed alphabet
of local constraint coefficients. The latter hardness statement retains an
explicitly nonconvex upper quadratic row. Exponential algebraic degree and
rational-output lower bounds explain why exact field output, numerical degree
and accuracy bits must be separated. The growing-leader W[1]/ETH
classification and the sparse-shadow construction have prior sources and are
attributed as such; the section does not claim every boundary as new.

\paragraph{Complete scalar computations.}
We provide exact convex and nonconvex scalar specializations, with original
response reconstruction and both upper selection conventions. A second solver
independently enumerates stationary faces in original coordinates and solves
the same continuous leader task, under an explicit nonsingularity restriction.
This is a validation baseline, not a new face-enumeration principle.
The experiments separate heterogeneous nonconvex instances from large repeated
populations, and report the preprocessing cost and unfavorable timings of
dense screening. These implementations illustrate specific specializations;
they do not implement the general real-algebraic algorithms.

Throughout, a polynomial bound for fixed structural dimensions can have an
exponent depending on those dimensions. No fixed-parameter tractability claim
is made. Polynomial degrees are numerical degrees, not the bit lengths of
sparse binary exponents. The precise models, selection conventions and
arithmetic output definitions follow below.

\subsection{Related work and the structural question}\label{subsec:related}

Classical positive results for optimistic linear bilevel optimization fix
the number of \emph{follower} variables. Liu and Spencer~\cite{LiuSpencer1995} give a
polynomial algorithm under that restriction; Deng~\cite{Deng1998}
discusses the resulting complexity theory. This restriction differs from
fixing the number of leader variables while allowing the follower dimension
to grow. Indeed, Sugishita and Carvalho~\cite{SugishitaCarvalhoIPCO2026,SugishitaCarvalho2026} prove
NP-completeness with one leader variable, unit-bounded variables, and no
additional upper constraints for a linear bilevel class with coupled
follower constraints. Their polynomial algorithm for finding a local
optimum does not give global tractability.

Ketkov and Prokopyev~\cite{KetkovProkopyev2026} distinguish fixed follower
variable and constraint counts, optimistic and pessimistic responses, and
linear and quadratic models. Their optimistic convex-quadratic positive
result fixes the follower dimension and imposes their stated convex
quadratic upper structure. Neither that result nor a restriction on the
total follower constraint count encompasses a model with arbitrarily many
local follower variables and local bounds, few shared constraints, and a
possibly nonconvex polynomial aggregate cost. In particular, counting
shared resource rows separately from local rows is essential. Their
pessimistic convention also illustrates why upper feasibility must be
specified before complexity statements are compared.

The local elimination tools have direct antecedents. Megiddo and
Tamir~\cite{MegiddoTamir1993} use low-dimensional multiplier arrangements
for separable convex quadratic optimization and for fixed-dimensional
quadratic blocks with shared constraints. The local quadratic formulas
and multiplier partitions used here build on this principle. Exact sign
determination and quantifier elimination in fixed dimension are classical
real algebraic algorithms~\cite{BasuPollackRoy1996}. The structural question
is how these ingredients provide a description of \emph{globally} optimal
responses when the aggregate cost is nonconvex, and how that description
supports the upper optimization. Stationarity alone cannot supply this
last step.

Global polynomial bilevel optimization also has direct antecedents beyond
KKT reformulations. Jeyakumar et al.~\cite{JeyakumarEtAl2016} and
Nie, Wang and Ye~\cite{NieWangYe2017} develop convergent semidefinite
relaxations under their respective assumptions. The latter use globally
quantified follower-value comparisons in a semi-infinite reformulation.
Nie et al.~\cite{NieEtAl2021} represent multipliers by polynomial or rational
expressions and combine them with feasible extensions and exchange steps.
Nie, Ye and Zhong~\cite{NieYeZhong2026} use sparse multiplier supports and
disjunctive decompositions for linearly constrained followers. Thus neither
global comparison, rational multiplier expressions nor conic support reduction
is introduced here. The distinction is quantitative: we eliminate the
\emph{primal} local variables as well, and bound the number of joint regimes
in a dimension independent of their total count. A support bound based on the
rank of the entire follower constraint matrix can still grow like the follower
dimension when that matrix contains all box normals. Our local support bound
uses the fixed block dimension. Global polynomial subproblem convergence and
a polynomial rational-bit bound for this compressed class are different
conclusions.

Fixed rank of the entire quadratic matrix is another established tractable
restriction for box optimization \cite{HladikCernyRada2021}. It differs from a
positive diagonal or block-diagonal matrix plus fixed-rank coupling, whose
total rank may grow with the follower count. Accordingly our result is not
tractability from low-rank nonconvexity alone: the local curvature and the
constraint structure are used in the proof.

For polynomial local costs, earlier work already supplies relevant
approximation methods. Hochbaum and Shanthikumar~\cite{HochbaumShanthikumar1990}
give logarithmic accuracy dependence for approximating the solution vector
of continuous separable allocation under their function-oracle assumptions,
with numerical dependence on a bound for constraint subdeterminants.
Vigneron~\cite{Vigneron2014} gives fully polynomial approximation schemes
for global optimization of sums of nonnegative algebraic functions of
constant description complexity, with polynomial dependence on inverse
relative error. A high-accuracy follower evaluation at a
fixed leader is different from a global approximation algorithm for the
induced, possibly signed upper objective. Both the numerical degree and
the accuracy dependence must be tracked in making that comparison.

Cost-near-optimal follower behavior is an established form of decision
uncertainty. Besan\c{c}on, Anjos, and Brotcorne~\cite{BesanconAnjosBrotcorne2024}
study robustness of upper feasibility against such responses;
Beck, Ljubi\'c, and Schmidt~\cite{BeckLjubicSchmidt2023} place this model in
the broader uncertainty literature. The robust model below also allows a
worst-case upper objective. The structural contribution concerns exact
compression of its adversarial problems, including nonconvex nominal
followers, rather than the introduction of near-optimality robustness.

Finally, polynomial-size optimality certificates need not produce
polynomial-time optimization. Buchheim~\cite{Buchheim2023} proves NP
membership for linear bilevel optimization and polynomial encoding bounds
for valid big-$M$ values in optimistic KKT reformulations. Practical
reformulations, branching, cuts, and decomposition are surveyed by
Kleinert et al.~\cite{KleinertEtAl2021}. They provide the computational
context for the specializations studied here. Polynomial time for fixed
structural dimensions is a worst-case statement; it is not by itself a
claim of favorable numerical performance.

Recent first-order work addresses several distinct guarantees. Chen, Ji and
Zhang~\cite{ChenJiZhang2026} study condition-number dependence for approximate
stationarity with a nonconvex upper problem and a strongly convex lower
problem. Wu et al.~\cite{WuEtAl2026} study stationarity with a uniformly
convex unconstrained follower; their Lemma~4.2 already gives the response
exponent $1/(p-1)$ for uniform convexity of order $p$. Our $1/P$ bound
extends this exponent to the structured polynomial model with moving affine
resources and a constant of polynomial encoding length
(\cref{prop:accuracy-response-modulus}). Global first-order guarantees also
exist: Xiao and Chen~\cite{XiaoChen2025} prove convergence of penalty-based
bilevel gradient descent under joint or blockwise Polyak--\L{}ojasiewicz
conditions on the penalty reformulation and accompanying assumptions,
with approximate lower optimality. Our global rational-bit guarantees use
fixed structural dimensions and do not require these landscape conditions.
Large populations have also motivated computational decomposition: Flocco,
Schiewe and Gabriel~\cite{FloccoEtAl2026} develop nested Benders decomposition
for many linear followers that decouple at a fixed leader. Their model and
computational objective differ from the jointly coupled polynomial follower
studied here.

Table~\ref{tab:prior-comparison} summarizes the closest comparisons.
To the best of our knowledge, the cited work does not establish the combined
exact quadratic-block or global accuracy-bit guarantees stated here with
unbounded follower dimension and the specified fixed shared dimensions.
This is a claim about those complete theorem classes and their outputs,
not priority for multiplier arrangements, global optimality comparisons,
polynomial inversion, convex conjugation or fixed-dimensional real algebra.

\begin{table}[t]
\centering\small
\caption{Relationship to selected prior work. The right column identifies the
additional guarantee under this paper's restrictions, not a claim that all
models in the cited work are subsumed.}
\label{tab:prior-comparison}
\begin{tabular}{@{}>{\raggedright\arraybackslash}p{.23\linewidth}>{\raggedright\arraybackslash}p{.32\linewidth}>{\raggedright\arraybackslash}p{.39\linewidth}@{}}
\toprule
Prior work & Established result or method & Distinction in this paper\\\midrule
Liu--Spencer; Ketkov--Prokopyev
\cite{LiuSpencer1995,KetkovProkopyev2026} &
Positive bilevel results with fixed follower dimension; further distinctions
by total constraint count and selection convention. &
Growing follower dimension and local row counts; fixed leader, local block,
shared resource and aggregate dimensions.\\[5pt]
Megiddo--Tamir; Nie--Ye--Zhong
\cite{MegiddoTamir1993,NieYeZhong2026} &
Low-dimensional quadratic multiplier arrangements; sparse multiplier
supports and KKT disjunctions. &
Polynomially many jointly realizable local regimes; primal recovery and
comparison of all global responses in fixed dimension.\\[5pt]
Hochbaum--Shanthikumar; Vigneron
\cite{HochbaumShanthikumar1990,Vigneron2014} &
High-accuracy separable allocation; approximation schemes for sums of
nonnegative algebraic functions. &
Global signed polynomial upper optimization in accuracy bits, fixed shared
coupling, and rational leader recovery with explicit feasibility guarantees.\\[5pt]
Besan\c{c}on--Anjos--Brotcorne
\cite{BesanconAnjosBrotcorne2024,BesanconAnjosBrotcorne2021} &
Cost-near-optimal response robustness and its complexity framework. &
Exact nonconvex measurement-fiber compression, with a fixed measurement count
separately for each criterion.\\[5pt]
Giesen et al.; Bemporad--Filippi
\cite{GiesenJaggiLaue2012,GiesenEtAl2012,BemporadFilippi2001} &
Certified approximate parametric paths and optimizer-error bounds. &
A saturation-based cost grid with a condition-dependent global upper bound
and no numerical incentive-magnitude factor beyond input encoding.\\[5pt]
Gardiner--Lucet; Moehle et al.
\cite{GardinerLucet2010,MoehleEtAl2023} &
Univariate piecewise quadratic convex-envelope constructions. &
An exact scalar bilevel specialization retaining original contacts, upper
feasibility, selection and attainment, with reproducible comparisons.\\
\bottomrule
\end{tabular}
\end{table}

\paragraph{Organization and notation.}
Section~\ref{sec:exact} proves exact compression and its extensions;
Section~\ref{sec:robust-screening} develops robustness and certified dense
recovery. Section~\ref{sec:accuracy} proves the global approximation results,
and Section~\ref{sec:boundaries} establishes their structural and arithmetic
limits. Section~\ref{sec:computation} gives complete scalar algorithms and
experiments. The appendices contain full supporting proofs for polyhedral
elimination, inverse approximation, quantitative resource bounds and path
geometry. The principal dimension parameters are $r$ (leader), $N$ (total
follower), $d$ (maximum local block), $s$ (aggregate) and $k$ (shared rows).
Only the parameters explicitly fixed in each theorem are held constant.
Symbols for matrices are local to their model; in particular $H$ denotes an
upper polynomial in Section~\ref{sec:accuracy}, whereas it denotes shared
inequality normals in the quadratic-block model.

\subsection{The quadratic-block model}\label{subsec:model}

Every polynomial datum in this model is supplied explicitly with rational
coefficients.

The leader chooses $x\in C\subset\R^r$, where $C$ is a compact set
specified by a rational quantifier-free polynomial formula. The follower
vector is $z=(z_1,\ldots,z_B)$, with $z_b\in\R^{d_b}$,
$d_b\le d$, and total dimension $N=\sum_{b=1}^B d_b$. Both $B$ and $N$
may grow. For the primary model, let
\begin{align}
 P_b(x)&=\{y\in\R^{d_b}:E_b y=e_b(x),\ G_b y\le h_b(x)\},
 \label{eq:local-polytope}\\
 P(x)&=\left\{z:z_b\in P_b(x)\ (b=1,\ldots,B),\
                     Az=b(x),\ Hz\le a(x)\right\}.
 \label{eq:follower-polytope}
\end{align}
The matrices $E_b,G_b,A,H$ are rational and independent of $x$.
There are $k_{=}$ shared equality rows and $k_{\le}$ shared inequality
rows, and $k=k_=+k_{\le}$. The number of local rows is unrestricted.
All right-hand sides are polynomial in $x$. Among the local inequalities
are coordinate bounds
$\ell_{bi}(x)\le z_{bi}\le u_{bi}(x)$, with polynomial endpoints
satisfying $\ell_{bi}(x)\le u_{bi}(x)$ on $C$.
Thus the union of the follower feasible sets is bounded.
A set $P_b(x)$ or $P(x)$ may be empty; follower feasibility is not a
promise at every leader.

Let $U(x)\in\Q[x]^{s\times N}$, and partition its columns into blocks
$U_b(x)$. For rational polynomial vectors $c_b(x)\in\Q[x]^{d_b}$,
the follower objective is
\begin{equation}\label{eq:quadratic-objective}
 f(x,z)=\sum_{b=1}^B\left(\frac12 z_b\trans Q_b(x)z_b+
                c_b(x)\trans z_b\right)+\phi(x,U(x)z).
\end{equation}
Each $Q_b(x)$ is a symmetric polynomial matrix positive definite for
every $x\in C$. The polynomial $\phi(x,w)$ can be nonconvex in $w$.
Its argument $w=U(x)z\in\R^s$ collects the shared aggregate quantities.
Positive definiteness is imposed on the local blocks, not on the entire
follower Hessian. In particular, the follower need not have a unique
minimizer.

The upper objective $F(x,z)$ and upper rows $g_j(x,z)\le0$
($j=1,\ldots,m$) are explicitly encoded rational polynomials.
Their monomials may involve many different follower coordinates; neither
$m$ nor their support sizes are structural parameters. Unless a theorem
says otherwise, the fixed parameters of the exact quadratic model are
\begin{equation}\label{eq:structural-parameters}
                    (r,d,s,k_=,k_{\le}).
\end{equation}
The scalar local model has $d=1$ and only interval constraints within a
block. Extensions with polynomially varying constraint normals will be
stated separately. The approximation results use strictly convex
univariate polynomial local costs and a convex aggregate cost, with their
own leader-domain and data assumptions; those assumptions are not
implicit extensions of \eqref{eq:quadratic-objective}.

\subsection{Follower selection and upper feasibility}\label{subsec:semantics}

For $P(x)\ne\varnothing$, compactness gives the value and solution set
\begin{equation}\label{eq:follower-value}
 v(x)=\min_{z\in P(x)} f(x,z),\qquad
 S(x)=\{z\in P(x):f(x,z)=v(x)\}.
\end{equation}
Put $S(x)=\varnothing$ if $P(x)=\varnothing$, and use $v(x)$ only
on feasible leaders. The follower always optimizes over $P(x)$ without
upper constraints. Define the optimistic feasible set and value by
\begin{align}
 \mathcal O&=\{(x,z):x\in C,\ z\in S(x),\
                             g_j(x,z)\le0\ (j=1,\ldots,m)\},
 \label{eq:optimistic-set}\\
 \operatorname{val}_{\rm opt}&=\inf_{(x,z)\in\mathcal O} F(x,z).
 \label{eq:optimistic-value}
\end{align}
Thus the leader may select a favorable global response that satisfies its
upper rows. A minimizer is asserted only when attainment is established.

For pessimistic optimization, upper feasibility is universal:
\begin{align}
 D_0&=\{x\in C:S(x)\ne\varnothing,\
             g_j(x,z)\le0\text{ for all }z\in S(x),\ j=1,\ldots,m\},
 \label{eq:pessimistic-domain}\\
 W_0(x)&=\max_{z\in S(x)}F(x,z),\qquad
 \operatorname{val}_{\rm pes}=\inf_{x\in D_0}W_0(x).
 \label{eq:pessimistic-value}
\end{align}
The maximum at a fixed feasible leader is attained. The infimum over
leaders need not be. An upper constraint cannot discard an inconvenient
follower optimum before $W_0$ is computed.

Given a nonnegative polynomial budget $\rho(x)$ on $C$, the set of
cost-near-optimal responses is
\begin{equation}\label{eq:nearoptimal-set}
 S_\rho(x)=\{z\in P(x):f(x,z)\le v(x)+\rho(x)\}
\end{equation}
at feasible leaders, and is empty otherwise. Replacing $S(x)$ by
$S_\rho(x)$ in \eqref{eq:pessimistic-domain}--\eqref{eq:pessimistic-value}
defines $D_\rho$, $W_\rho$, and $\operatorname{val}_\rho$.
A bounded budget may itself be a leader coordinate when the leader model
permits it. Near-optimal points need not satisfy follower stationarity.
The robust compression theorem will restrict the number of affine
measurements of $z$ used by each upper criterion; this is an additional
hypothesis, distinct from the unrestricted polynomial upper data in the
exact optimistic model.

For revenue maximization, the corresponding optimistic objective is a
supremum and the pessimistic objective maximizes the minimum revenue over
all responses. Equivalently, negate revenue in the minimization
conventions above. The exact algorithms report infeasibility separately
and state whether the finite optimum value is attained. Approximation
algorithms have the guarantees stated in their respective theorems:
an original-infeasibility certificate, output satisfying relaxed upper
rows, and an empty inner surrogate have different meanings. In particular,
neither relaxed output nor an empty inner surrogate decides original
feasibility in general.

\subsection{Input length, degree, and output}\label{subsec:encoding}

The bit model counts arithmetic on binary-encoded rational data. Let $L$
be the total rational input encoding length, including dimensions, row
indices, monomial lists, exponents, and the Boolean formula for $C$.
Let $\delta\ge2$ bound the actual total degree of every input polynomial.
The exact algorithms have running time polynomial in $(L,\delta)$ for
fixed \eqref{eq:structural-parameters}. A claim of polynomial time in $L$
alone assumes that $\delta$ is polynomially bounded by $L$, as for dense
or unary-exponent encodings. A sparse monomial representation with
unrestricted binary exponents does not have that property. Polynomial
local costs in the approximation results are supplied densely; a
power-specific result will state its dependence on the numerical power.

A common-field algebraic output consists of a square-free polynomial
$p\in\Z[T]$, a rational isolating interval selecting a real root
$\alpha$, and rational polynomials $p_i,q_i$ such that each returned
coordinate equals $p_i(\alpha)/q_i(\alpha)$ with $q_i(\alpha)\ne0$.
The guarantee includes the degree of $p$, all coefficient bit lengths,
and the combined encoding of all coordinates. They are bounded
polynomially in $(L,\delta)$ in the stated fixed dimensions, and hence
polynomially in $L$ under the degree-encoding condition above. The
polynomial $p$ need not be a minimal polynomial. It suffices that its
chosen root specifies the same real extension for every recovered
coordinate. If several adversarial problems are requested independently,
a common-field guarantee for each leader--witness pair does not imply a
polynomial-degree field containing all their unrelated witnesses at once.

An accuracy-bit guarantee takes a positive rational tolerance $\eps$,
whose encoding contributes to the running time. Writing $\eps=2^{-\beta}$
illustrates the distinction between a bound polynomial in $\beta$ and
one polynomial in $1/\eps$. In the base approximation model, the output
is a rational leader that is exactly feasible for the base leader domain
and admits a feasible follower, with an additive guarantee for the true
induced upper objective. An approximate follower vector does not replace
the true response in that guarantee. Exact satisfaction of
response-dependent upper rows requires additional conditions; without
them, a theorem must specify its allowed upper violation. Exact
algebraic output, rational near-optimal output, and numerical evaluation
are different computational tasks.

The exponents of the polynomial bounds depend on the structural
parameters. For the decision versions this gives XP-type bounds,
$L^{h(r,d,s,k)}$, under the degree convention. It does not establish a
fixed-parameter bound $h(r,d,s,k)L^c$ with an absolute exponent $c$.

\subsection{Algebraic and polyhedral tools}\label{subsec:tools}

We use the following standard consequences of fixed-dimensional real
algebraic computation~\cite[Theorem 1.3.1 and Section 3.1.3]{BasuPollackRoy1996}.
For a rational polynomial family in a fixed total number of real
variables, realizable sign conditions, including zero signs, can be
enumerated in time polynomial in the number of polynomials, their degrees,
and coefficient bit lengths. A first-order formula with a fixed total
number of free and bound variables can be replaced by an equivalent
quantifier-free formula within the same kind of bound. Nonempty
semialgebraic sets admit algebraic samples of polynomial degree and
encoding length. Applied simultaneously to all free coordinates, sampling
gives the common-field representation described above. The total number
of variable occurrences introduced by quantified copies must be controlled;
fixing only the number of free variables is insufficient.

\begin{samepage}
One elementary encoding fact will be used repeatedly.
\begin{lemma}[Substitution in fixed dimension]\label{lem:substitution}
Let $v\in\R^q$, with $q$ fixed, and suppose that $Z_1(v),\ldots,Z_N(v)$
and $D(v)$ have explicit polynomial encodings of polynomial size and degree.
On a region where $D>0$, put $z_i=Z_i/D$. Let $p(x,z)$ be an explicitly
listed rational polynomial of total degree at most $\delta$, with the
coordinates of $x$ among those of $v$. Then
\[
             D(v)^\delta p(x,Z(v)/D(v))
\]
is a polynomial that can be constructed in time polynomial in the input
encoding lengths, the degrees of $Z_i,D$, and $\delta$. Its sign agrees
with that of $p(x,z)$ on the region.
\end{lemma}
\end{samepage}
\begin{proof}
An input monomial $a x^\nu z^\gamma$ becomes
$a x^\nu\prod_iZ_i^{\gamma_i}D^{\delta-|\gamma|}$.
Its degree is at most $\delta(1+T)$ if the degrees of $Z_i,D$ are at
most $T$. A polynomial of degree at most $\delta(1+T)$ in $q$ variables
has at most $\binom{\delta(1+T)+q}{q}$ monomials. This number is
polynomial for fixed $q$. Expanding and adding the explicit input
monomials therefore takes polynomially many rational operations;
coefficient bit lengths grow at most polynomially under these products
and sums. Positivity of $D^\delta$ proves sign preservation.
\end{proof}

For polyhedral limits we use Hoffman's bound~\cite{Hoffman1952}.
Given fixed matrices $M,T$, there is a constant $H_{M,T}$ such that,
whenever $K(a,b)=\{y:My\le a,\ Ty=b\}$ is nonempty,
\begin{equation}\label{eq:hoffman}
 \dist(y,K(a,b))\le H_{M,T}
          \bigl(\norm{(My-a)_+}+\norm{Ty-b}\bigr).
\end{equation}
Any fixed choices of finite-dimensional norms can be used, with a changed
constant. The constant is uniform in the right-hand sides. Equalities
can be encoded as opposite inequalities. In the exact attainment proof,
only existence of this constant for each instance is needed. The
approximation proofs require quantitative encoding bounds and will
establish them explicitly. Neither form of the argument grants a uniform
constant when the constraint normals vary with the leader.
